\section{Weighted height data and convex subobjects}
\label{sec:height-data}

This section packages the layer counts by original colors and proves
the two properties needed for their later comparisons.  When the
boundary layers vanish above a chosen level, we can choose a
representative in the spherical coset with the same upper-layer
vanishing on its type.  A uniform lower bound on the remaining layer then
produces actual subrepresentations through a convex-hull argument.

For $x\in A$ and a simple object $L$ of $\mathscr C$, consider the
left projective complex
\begin{equation}
 E_L(x)=B_xP_{o,L},\qquad Y_L^d(x)=Y^d(E_L(x)).
 \label{eq:height-data-frame-family}
\end{equation}
The family member with label $L$ is assigned weight
$w_L=\dim_{\mathscr C}L>0$.  All layer multiplicities below are
computed in minimal models.  Their independence of those models
follows from Lemma~\ref{lem:layer-structure}.

\subsection{Label tensoring and proportionality}

\begin{lemma}[Tensoring and canonical truncations]
\label{lem:height-data-tensor-truncation}
For $H\in\mathscr C$, right tensoring of all projective labels by
$H$ defines a $t$-exact functor $T_H$ for the layer $t$-structure.
It commutes naturally with its truncations, layers, and all the
color twist functors and their inverses.

In particular, if a subobject is constructed by applying a word
of twists, taking a layer truncation, and applying the inverse
word, this construction commutes with $T_H$ as a morphism into
the ambient object.  No choice of minimal model or Gaussian
elimination is required to be natural.
\end{lemma}

\begin{proof}
Proposition~\ref{prop:twist-action} supplies the additive
functor $P_{s,X}\mapsto P_{s,X\otimes H}$ and its natural
commutation with twists.  It extends termwise to complexes and
commutes with cones and shifts.  A positive-length differential
entry remains of positive length after tensoring and decomposing
the labels.  Hence a minimal model stays minimal, with the same
cochain and internal placements.  Both halves of the layer
$t$-structure are preserved, proving $t$-exactness.

The images under a $t$-exact functor of the two objects in a
truncation triangle still have the required bounds.  The
uniqueness of $t$-truncations therefore gives natural
commutation with $\tau_{\leq d}$ and $\tau_{\geq d}$, and
with the induced layers.  Commutation with inverse twists
follows from the natural commutation with the positive
equivalences.  Finally a truncation-defined subobject is a
composite of a canonical truncation map with the unit or
counit of the inverse equivalences.  The preceding natural
isomorphisms commute with that composite.  One may subsequently
represent the resulting pure objects and subobjects by their
coefficient arrays using Lemma~\ref{lem:layer-structure}.
\end{proof}

\begin{lemma}[Proportionality of the frame family]
\label{lem:family-proportionality}
There are nonnegative vectors
$m^d(x)\in\R^{\widehat S}$ such that the weighted generator
dimension vector of the layer family is
\begin{equation}
 \sum_L w_L\dim Y_L^d(x)=\phi m^d(x).
 \label{eq:height-data-proportionality}
\end{equation}
The weighted top and socle dimension vectors of these layers
also have the form $\phi v$.  Consequently, if a top or socle
occurs anywhere in a color in the family, then at every unfolded
vertex of that color it occurs in at least one member of the
family.

For each $x$ the vectors $m^d(x)$ have finite support in $d$.
Uniformly over $x$, $d$, and the colors, each coordinate has
only finitely many possible values in a bounded interval.
Operations on the left by $\sigma_i^{\pm1}$ change $m^d$ only
in coordinate $i$, and
\begin{equation}
 m_o^d(x)=
 \begin{cases}1,&d=0,\\0,&d\ne0.\end{cases}
 \label{eq:height-data-frame-coordinate}
\end{equation}
On folded signed counts, either twist in color $i$ acts by the
coordinate reflection for $C$:
\begin{equation}
 \sum_d(-1)^d m^d(\sigma_i^{\pm1}x)
   =s_i\left(\sum_d(-1)^d m^d(x)\right).
 \label{eq:height-data-signed-reflection}
\end{equation}
\end{lemma}

\begin{proof}
Let $\mathbf1$ be the unit label.  Natural label equivariance
gives
\[
 E_L(x)\simeq T_L E_{\mathbf1}(x),\qquad
 Y_L^d(x)\simeq T_LY_{\mathbf1}^d(x).
\]
For an object $X\in\mathscr C$ and a simple label $H$, rigidity
and semisimplicity give
\begin{equation}
 \sum_L w_L[X\otimes L:H]
       =w_H\dim_{\mathscr C}X.
 \label{eq:height-data-weighted-fusion}
\end{equation}
Indeed $[X\otimes L:H]$ is the multiplicity of $L$ in
$X^*\otimes H$, by adjunction.  Summing the dimensions of these
simple summands gives
$\dim(X^*\otimes H)=\dim(X)\dim(H)$.

In a minimal representative of $E_{\mathbf1}(x)$, collect the
projective labels at each color and pair of placements into an
object $X\in\mathscr C$.  Tensoring that summand by $L$ gives
$P_{s,X\otimes L}$ with unchanged placements.  Applying
\eqref{eq:height-data-weighted-fusion} and summing over the
placements of a fixed level gives
\eqref{eq:height-data-proportionality}.

The same argument applies to tops and socles, as follows.  For
any such projective complex and any color $s$, its test at
$(s,H)$ is obtained by applying $\Hom_{\mathscr C}(H,-)$ to a
complex of objects of $\mathscr C$: a term of type $P_{t,X}$
contributes $D^a_{st}\otimes X$ at path length $a$, with the
given placements.  The maps are induced by the same
$D$-multiplications as the projective differential.  Tensoring
labels by $L$ tensors this entire label complex on the right
by $L$.  Tensoring in the semisimple category is exact, so it
commutes with its cohomology.  Apply
\eqref{eq:height-data-weighted-fusion} to the cohomology in the
two extremal total degrees from
Lemma~\ref{lem:layers-extremal-tests}.  This proves the
proportionality for both tops and socles.  Since every $w_H$
and every family weight is positive, a nonzero proportional
coordinate at one label is nonzero at all labels of that
color, proving the occurrence assertion.

There are finitely many family members and finitely many
generators in each of their complexes, proving finite level
support.  Fixing one simple label $H$ in
\eqref{eq:height-data-proportionality} writes $m_s^d$ as a
nonnegative integer combination of the finitely many positive
numbers $w_L/w_H$.  Below a fixed bound, each integer
coefficient is bounded.  Thus there are only finitely many
values; irrational label dimensions cause no accumulation
through cancellation.

The outside-color statement follows from the right-simple
test in Lemma~\ref{lem:layers-inverse-amplitude}.  In
particular $U_{o,H}B_x\simeq U_{o,H}$, because the frame is
never acted on.  Its tensor product with $P_{o,L}$ is one
copy of $\C$ in placements $(0,0)$ when $H=L$, and is zero
otherwise.  The weighted count is therefore $w_H$ in level
zero and zero in the other levels.  This proves
\eqref{eq:height-data-frame-coordinate}.

Finally, the signed calculation in
Lemma~\ref{lem:positive-layer} gives the unfolded reflection
on generator counts.  Its inverse acts by the same reflection,
which is an involution.  The identity $\widetilde C\phi=\phi C$ in
\eqref{eq:category-intertwining} folds this to
\eqref{eq:height-data-signed-reflection}.
\end{proof}

For later normalizations set
\begin{equation}
 D_0=\sum_Lw_L^2.
 \label{eq:height-data-D0}
\end{equation}
In particular the ordinary dot product of two proportional
arrays is $(\phi p)^{\mathsf t}(\phi y)=D_0p^{\mathsf t}y$.

\subsection{Caps and spherical accessibility}

\begin{definition}[Cap]
\label{def:height-data-cap}
For $q\in\Zbb$ and $K\subseteq\widehat S$, the data at a
representative $x$ are \emph{capped at $q$ on $K$} if
\[
 m_i^d(x)=0\qquad(i\in K, d>q).
\]
For a set of colors $J$, let $\partial J$ denote its outside
neighbors in the framed diagram; it includes the frame when
the frame is a neighbor.  Positivity of the family weights
makes a cap equivalent to the absence of generators in all
the corresponding unfolded vertices and levels in every
member of the family.
\end{definition}

The exterior coordinates of every $m^d$ are constant on a
coset $A_Jx$, by Lemma~\ref{lem:family-proportionality}.
All caps imposed on sets outside $J$ are therefore independent
of the choice of representative in that coset.

We first isolate how a cap interacts with layers outside the
operated region.  A complex supported on colors outside
$J\cup\partial J$ has zero test at every unfolded vertex
of $J$.  Its positive evaluation and negative coevaluation
terms in $J$ vanish, so every twist in $J$ acts on that
complex as the identity.

\begin{lemma}[A layer under a capped negative operation]
\label{lem:height-data-capped-layer}
Let $E$ be a bounded minimal projective complex with no
generators above $q$ on $J\cup\partial J$.  If $G$ is a
word of negative twists in $J$, then
\begin{equation}
 Y^q(GE)\simeq Y^q\bigl(G(Y^q(E))\bigr).
 \label{eq:height-data-capped-layer}
\end{equation}
Moreover $GE$ still has no generators above $q$ on
$J\cup\partial J$.
\end{lemma}

\begin{proof}
Let $R=\tau_{\geq q+1}E$.  By the cap, every generator of
$R$ is outside $J\cup\partial J$, so $GR\simeq R$ by the
preceding observation.  Negative twists do not increase
levels.  Apply $G$ to
\[
 \tau_{\leq q}E\longrightarrow E\longrightarrow R.
\]
The first term has levels at most $q$, and the third is the
unchanged complex $R$ in levels at least $q+1$.  The long
exact sequence identifies level $q$ of $GE$ with level $q$
of $G\tau_{\leq q}E$, and all higher layers with the
corresponding layers of $R$.  These higher layers remain
outside the operated region, proving cap preservation.
Now apply $G$ to
\[
 \tau_{\leq q-1}E\longrightarrow\tau_{\leq q}E
       \longrightarrow Y^q(E).
\]
Its first term still has levels at most $q-1$, giving
\eqref{eq:height-data-capped-layer} by the long exact
sequence.  This argument does not require that $G$ be a
reduced word.
\end{proof}

\begin{lemma}[Cap accessibility]\label{lem:cap-accessibility}
Let $J\subseteq S$ be spherical, let $q\in\Zbb$, and suppose
the data at $x$ are capped at $q$ on $\partial J$.  There
exists $z\in A_Jx$ whose data are capped at $q$ on
$J\cup\partial J$.

Every negative generator in $J$ preserves an existing cap
at $q$ on $J\cup\partial J$.  At a representative with
this joint cap, a positive $F_i$, $i\in J$, also preserves
it whenever
\[
 \operatorname{top}_\alpha Y_L^q=0
 \qquad\text{for every $L$ and every $\alpha$ of color $i$}.
\]
\end{lemma}

\begin{proof}
Negative preservation follows by applying
Lemma~\ref{lem:height-data-capped-layer} to each family
member.  For the positive assertion use the same high
quotient $R=\tau_{\geq q+1}E_L$.  It is unchanged by
$F_i$.  On $\tau_{\leq q}E_L$ the only possible level
$q+1$ contribution is the upper piece of $F_iY_L^q$:
input levels below $q$ can rise by at most one.  By
Lemma~\ref{lem:positive-layer}, this upper piece is the
old top at color $i$, and vanishes under the hypothesis.
The low part therefore stays at levels at most $q$, while
the high quotient stays outside $J\cup\partial J$.
This proves the positive cap assertion.

For accessibility there is nothing to prove if $J$ is
empty or already capped.  Otherwise let $t>q$ be the
largest level with nonzero data on $J$.  The joint cap at
$t$ holds on $J\cup\partial J$.  In each $Y_L^t$ the
part supported on $\widetilde J$ is a direct summand,
because all boundary generators at that level vanish.
If these parts are nonzero, some member has a nonzero
top at a color $i\in J$.  Apply $F_i^{-1}$ and repeat
while level $t$ persists.  All these moves preserve the
joint cap at $t$.

We show that the chosen colors form a reduced expression
in $W_J$, read in their order of action.  Suppose the
current negative word represents a reduced Coxeter
expression $w$.  Every left descent $j$ of $w$ can occur
as its last-action letter, since reduced expressions are
related by braid moves.  The corresponding negative
Artin words are related by the same braid moves.  In
that expression, let $E'_L$ be the family immediately
before its last $F_j^{-1}$.  All preceding negative
moves preserve the joint $t$-cap.  For $\alpha$ in color
$j$, Formula~\eqref{eq:layers-negative-test-shift} and
Lemma~\ref{lem:layers-extremal-tests} give
\[
 H^{t+2}(e_\alpha F_j^{-1}E'_L)
       =H^{t+3}(e_\alpha E'_L)=0.
\]
Thus the current level-$t$ top vanishes in every descent
color.  A chosen nonzero top must be at a left ascent,
so appending its action increases Coxeter length.
This proves the induction from the empty word.

The finite group $W_J$ has a longest length.  Therefore
level $t$ must disappear after at most that many chosen
moves.  Negative operations cannot recreate a larger
level.  Repeating at the next largest excess level
terminates, since the initial largest level is finite
and each stage decreases it toward $q$.  All operations
lie in $A_J$, and the exterior data stay fixed, proving
the required accessibility.
\end{proof}

\subsection{Convex subobjects of a capped layer}

Let $J$ be spherical.  Given nonnegative data on its
complement, there is a unique vector agreeing with those
data outside $J$ and harmonic on $J$, namely satisfying
$(Cp)_J=0$.  The nonnegative inverse $C_J^{-1}$ gives
nonnegative interpolation, as in
Lemma~\ref{lem:cartan-positive-null}.  The reflection
action of $W_J$ on the full coordinate space fixes such
a vector $p$.

The use of reflection words to select subobjects has a parallel
in the torsion-pair description of preprojective-module polytopes
\cite[Sections~3.2, 5.4 and~5.6]{BaumannKamnitzerTingley2014}.
Here we construct the required graded subrepresentations by
transporting layer truncations through the twists.

\begin{proposition}[Convex subobjects]
\label{prop:convex-subobjects}
Let $J\subseteq S$ be spherical, and suppose that at
$x$ the data are capped at $q$ on $J\cup\partial J$.
Let $p$ agree with $m^q(x)$ outside $J$ and be harmonic
on $J$, and put $e=m^q(x)-p$, which is supported on $J$.
Assume that for \emph{every} representative $z\in A_Jx$
whose data are capped at $q$ on $J\cup\partial J$,
\begin{equation}
 m_j^q(z)\geq p_j\qquad(j\in J).
 \label{eq:height-data-universal-cap-inequality}
\end{equation}
Then $e_J$ belongs to the convex hull in $\R^J$ of
vectors $s_J$ with the following property: $s$ is zero
outside $J$, and there are graded subrepresentations
\begin{equation}
 \mathcal S_L\subseteq Y_L^q(x),\qquad
 \supp\mathcal S_L\subseteq\widetilde J,
 \qquad
 \sum_Lw_L\dim\mathcal S_L=\phi s.
 \label{eq:height-data-convex-subrepresentations}
\end{equation}
One may use a finite collection of such vectors,
together with the zero vector.

If all inequalities in
\eqref{eq:height-data-universal-cap-inequality} are
strict, $e_J$ is in the ordinary interior in $\R^J$
of such a convex hull.  For $J=\varnothing$ the
interior convention is that $\{0\}$ is open in
$\R^0$.
\end{proposition}

\begin{proof}
The exterior coordinates of $m^q$ are fixed on $A_Jx$,
so the vector $p$ in the hypothesis is the same for
every representative under consideration.  For each
$w\in W_J$, choose a reduced expression and let $G_w$
be its negative lift acting on the left.  Its action
on signed folded counts is $w$, because each negative
letter acts by the same coordinate reflection as its
positive inverse.  Both $G_w$ and its inverse have
the reduced-word amplitude from
Corollary~\ref{cor:reduced-amplitude}.

Put $M_L=Y_L^q(x)$ and $N_L=G_wM_L$.  The latter has
levels $q-1,q$.  Its truncation triangle is
\begin{equation}
 N_{L,\mathrm{low}}\longrightarrow N_L
       \longrightarrow N_{L,\mathrm{high}},
 \label{eq:height-data-two-layer-triangle}
\end{equation}
where the two end terms are pure at $q-1$ and $q$,
respectively.  The low term is supported on
$\widetilde J$: all outside-color multiplicities in
each layer are unchanged under $G_w$, and the input
was pure at $q$.

Apply $G_w^{-1}$ to
\eqref{eq:height-data-two-layer-triangle}.  Write its
end terms as $A_L$ and $C_L$, so that the triangle is
\[
 A_L\longrightarrow M_L\longrightarrow C_L.
\]
The amplitude bounds place $A_L$ in $[q-1,q]$ and
$C_L$ in $[q,q+1]$.  The long exact sequence at $q-1$
forces $H^{q-1}_{\mathrm{lay}}(A_L)=0$, since the
adjacent middle and right groups there vanish.  At
$q+1$ it forces
$H^{q+1}_{\mathrm{lay}}(C_L)=0$.  Thus both end terms
are pure at $q$, and the sequence at $q$ is short
exact:
\begin{equation}
 0\longrightarrow A_L\longrightarrow M_L
       \longrightarrow C_L\longrightarrow0
 \label{eq:height-data-subobject-sequence}
\end{equation}
in the shifted heart.  By
Lemma~\ref{lem:layer-structure}, $A_L$ is an actual
graded subrepresentation $\mathcal S_L$ of $M_L$.
Its support is still in $\widetilde J$, because a
twist at a color in $J$ sends a complex of projectives
of colors in $J$ to another such complex.

The subobject morphism in
\eqref{eq:height-data-subobject-sequence} is
\[
 G_w^{-1}\tau_{\leq q-1}(G_wM_L)
       \longrightarrow M_L.
\]
Lemma~\ref{lem:height-data-tensor-truncation} shows
that it is obtained from the unit-label morphism by
right tensoring with $L$, up to an isomorphism of
inclusions.  Applying the weighted fusion identity to
its generator spaces therefore gives
\eqref{eq:height-data-convex-subrepresentations} for
a vector $s=s_w$.  This argument uses canonical
truncations; it makes no choice of elimination
splittings simultaneously in the different labels.

The quotient family $C_L$ has weighted dimension
$\phi(m^q(x)-s_w)$.  Applying $G_w$ to it recovers
$N_{L,\mathrm{high}}$, which is pure at $q$.
Signed generator counts consequently give the folded
dimension of this high family as
\begin{equation}
 w\bigl(m^q(x)-s_w\bigr).
 \label{eq:height-data-reflected-quotient}
\end{equation}
There is no contribution here from higher layers of
the original $E_L(x)$ outside the operated region:
Lemma~\ref{lem:height-data-capped-layer} identifies
its level-$q$ family after $G_w$ with the level-$q$
family of $G_wM_L$, exactly the high family used in
\eqref{eq:height-data-reflected-quotient}.

The representative obtained by $G_w$ belongs to
$A_Jx$ and retains the joint cap, by
Lemma~\ref{lem:cap-accessibility}.  Applying the
universal hypothesis to
\eqref{eq:height-data-reflected-quotient}, and using
$wp=p$, gives
\begin{equation}
 \bigl(w(e-s_w)\bigr)_j\geq0\qquad(j\in J).
 \label{eq:height-data-chamber-inequalities}
\end{equation}
In the strict case every one of these coordinates is
strictly positive.

We finish in the finite-dimensional real vector
space of root coordinates on $J$.  The finite
Coxeter chamber theorem says that every linear
functional $f$ on this space can be written
\[
 f(v)=\sum_{j\in J}\lambda_j(wv)_j
 \quad\text{for some }w\in W_J\text{ and }
 \lambda_j\geq0.
\]
Indeed the nonnegative coordinate functionals form
the closed fundamental chamber in the dual
representation, whose $W_J$-translates cover
the dual space.  For the $w$ chosen for $f$,
\eqref{eq:height-data-chamber-inequalities} implies
\begin{equation}
 f(e)\geq f(s_w)\geq\min_{v\in\mathcal V} f(v),
 \qquad
 \mathcal V=\{s_u:u\in W_J\}\cup\{0\}.
 \label{eq:height-data-support-inequalities}
\end{equation}
The set $\mathcal V$ is finite.  If $e$ lay outside
its convex hull, a separating linear functional
would contradict
\eqref{eq:height-data-support-inequalities}.
Thus $e_J$ lies in that convex hull.

If the original inequalities were all strict, then
for nonzero $f$ at least one $\lambda_j$ is positive,
so $f(e)>f(s_w)\geq\min_{v\in\mathcal V}f(v)$.
A point in a finite convex hull which satisfies
this strict inequality for every nonzero functional
lies in its ordinary ambient interior: a boundary
point, or a convex hull in a proper affine
subspace, has a nonzero supporting functional
with equality.  This proves the strict assertion.
For $J=\varnothing$ the construction gives
$e=s_w=0$ and the stated zero-dimensional conclusion
directly.
\end{proof}
